\unitchapter{Paper 1 · Unit 1}{Teaching Aptitude}{p1-01}

\textbf{Expected questions: \textasciitilde5 (10 marks) · Target: 5/5}

\begin{sylbox}

\begin{itemize}
\tightlist
\item[$\square$]
  Teaching: concept, objectives, levels (memory, understanding, reflective), characteristics \& basic requirements
\item[$\square$]
  Learner\textquotesingle s characteristics (adolescent \& adult; academic, social, emotional, cognitive), individual differences
\item[$\square$]
  Factors affecting teaching (teacher, learner, support material, instructional facilities, environment, institution)
\item[$\square$]
  Methods of teaching: teacher-centred vs learner-centred; offline vs online (SWAYAM, SWAYAMPRABHA, MOOCs)
\item[$\square$]
  Teaching support system: traditional, modern, ICT-based
\item[$\square$]
  Evaluation systems: elements \& types; evaluation in CBCS; computer-based testing; innovations
\end{itemize}

\end{sylbox}

\needdiagram{figures/p1-01-00}{90pt}

\hypertarget{p1-01--1-what-is-teaching}{%
\section{1. What is Teaching?}\label{p1-01--1-what-is-teaching}}

Teaching is a \textbf{purposeful, interactive process} that brings a desirable change in the learner\textquotesingle s behaviour
(knowledge, skills, attitudes).

\diagramhere{figures/p1-01-00}

\begingroup\small

\begin{longtable}[]{@{}
  >{\raggedright\arraybackslash}p{(\columnwidth - 4\tabcolsep) * \real{0.1244}}
  >{\raggedright\arraybackslash}p{(\columnwidth - 4\tabcolsep) * \real{0.3175}}
  >{\raggedright\arraybackslash}p{(\columnwidth - 4\tabcolsep) * \real{0.1791}}@{}}
\toprule\noalign{}
\begin{minipage}[b]{\linewidth}\raggedright
\textbf{Variable type}
\end{minipage} & \begin{minipage}[b]{\linewidth}\raggedright
\textbf{What it is}
\end{minipage} & \begin{minipage}[b]{\linewidth}\raggedright
\textbf{Example}
\end{minipage} \\
\midrule\noalign{}
\endhead
\bottomrule\noalign{}
\endlastfoot
Independent & Teacher (plans, organises, leads) & Teacher\textquotesingle s planning \\
Dependent & Student (acts as per teacher) & Learner\textquotesingle s performance \\
Intervening & Content, method, A/V aids, environment & Smart board \\
\end{longtable}

\endgroup

\needdiagram{figures/p1-01-01}{55pt}

\hypertarget{p1-01--phases-of-teaching-philip-w-jackson}{%
\subsection{Phases of Teaching (Philip W. Jackson)}\label{p1-01--phases-of-teaching-philip-w-jackson}}

\diagramhere{figures/p1-01-01}

\begin{itemize}
\tightlist
\item
  \textbf{Pre-active}: fixing goals, choosing content, deciding method.
\item
  \textbf{Inter-active}: perception of class size, diagnosis, action (questioning, reinforcement).
\item
  \textbf{Post-active}: test design, measuring changes, changing strategies.
\end{itemize}

\needspace{179pt}

\hypertarget{p1-01--levels-of-teaching-most-asked}{%
\subsection{Levels of Teaching (most asked!)}\label{p1-01--levels-of-teaching-most-asked}}

\begin{asciiblock}{124pt}
\begin{Verbatim}[fontsize=\fontsize{8.2}{10.0}\selectfont]
                 ┌────────────────────────────────────────────┐
  Level 3        │ REFLECTIVE level — Hunt (1971)             │  Student-centred, problem solving,
  (highest)      │ "thoughtful" — critical, creative thinking │  highest insight
                 ├────────────────────────────────────────────┤
  Level 2        │ UNDERSTANDING level — Morrison (1931)      │  Memory + insight; "memory with
                 │ Generalisation, principles                 │  understanding"; teacher-centred
                 ├────────────────────────────────────────────┤
  Level 1        │ MEMORY level — Herbart (Herbartian steps)  │  Rote learning, thoughtless,
  (lowest)       │ Recall, recognition                        │  teacher-dominated
                 └────────────────────────────────────────────┘
\end{Verbatim}
\end{asciiblock}

\begingroup\small

\begin{longtable}[]{@{}
  >{\raggedright\arraybackslash}p{(\columnwidth - 8\tabcolsep) * \real{0.1237}}
  >{\raggedright\arraybackslash}p{(\columnwidth - 8\tabcolsep) * \real{0.1015}}
  >{\raggedright\arraybackslash}p{(\columnwidth - 8\tabcolsep) * \real{0.1861}}
  >{\raggedright\arraybackslash}p{(\columnwidth - 8\tabcolsep) * \real{0.1228}}
  >{\raggedright\arraybackslash}p{(\columnwidth - 8\tabcolsep) * \real{0.2113}}@{}}
\toprule\noalign{}
\begin{minipage}[b]{\linewidth}\raggedright
\textbf{Level}
\end{minipage} & \begin{minipage}[b]{\linewidth}\raggedright
\textbf{Proponent}
\end{minipage} & \begin{minipage}[b]{\linewidth}\raggedright
\textbf{Teacher role}
\end{minipage} & \begin{minipage}[b]{\linewidth}\raggedright
\textbf{Student role}
\end{minipage} & \begin{minipage}[b]{\linewidth}\raggedright
\textbf{Evaluation}
\end{minipage} \\
\midrule\noalign{}
\endhead
\bottomrule\noalign{}
\endlastfoot
Memory & Herbart & Dominant & Passive & Oral/\allowbreak{}written recall tests \\
Understanding & Morrison & Dominant & Semi-active & Objective + essay \\
Reflective & Hunt & Democratic/\allowbreak{}secondary & Active & Essay, problem-solving \\
\end{longtable}

\endgroup

Herbartian 5 steps: \textbf{Preparation → Presentation → Comparison/Association → Generalisation → Application}.
Morrison\textquotesingle s steps: \textbf{Exploration → Presentation → Assimilation → Organisation → Recitation}.

\needspace{144pt}

\hypertarget{p1-01--blooms-taxonomy-cognitive-domain-revised-2001-by-anderson--krathwohl}{%
\subsection{Bloom\textquotesingle s Taxonomy (Cognitive domain, revised 2001 by Anderson \& Krathwohl)}\label{p1-01--blooms-taxonomy-cognitive-domain-revised-2001-by-anderson--krathwohl}}

\begin{asciiblock}{89pt}
\begin{Verbatim}[fontsize=\fontsize{8.8}{10.7}\selectfont]
                        ▲  Creating      (design, construct, produce)
                       ▲▲  Evaluating    (judge, critique, justify)
                      ▲▲▲  Analysing     (compare, differentiate)
                     ▲▲▲▲  Applying      (use, execute, solve)
                    ▲▲▲▲▲  Understanding (explain, classify)
                   ▲▲▲▲▲▲  Remembering   (recall, list, define)
\end{Verbatim}
\end{asciiblock}

Original (1956): Knowledge → Comprehension → Application → Analysis → \textbf{Synthesis → Evaluation}.
Revised: Evaluation comes \textbf{before} Creating (Synthesis renamed Creating; nouns became verbs).

Three domains: \textbf{Cognitive} (Bloom), \textbf{Affective} (Krathwohl: Receiving → Responding → Valuing → Organisation → Characterisation),
\textbf{Psychomotor} (Simpson / Dave: Imitation → Manipulation → Precision → Articulation → Naturalisation).

\needspace{100pt}

\hypertarget{p1-01--maxims-of-teaching}{%
\subsection{Maxims of Teaching}\label{p1-01--maxims-of-teaching}}

\begin{itemize}
\tightlist
\item
  Known → Unknown · Simple → Complex · Concrete → Abstract · Particular → General · Whole → Part
\item
  Analysis → Synthesis · Empirical → Rational · Psychological → Logical · Induction → Deduction
\end{itemize}

\needdiagram{figures/p1-01-02}{55pt}

\hypertarget{p1-01--2-learners-characteristics}{%
\section{2. Learner\textquotesingle s Characteristics}\label{p1-01--2-learners-characteristics}}

\diagramhere{figures/p1-01-02}

\begin{itemize}
\tightlist
\item
  \textbf{Pedagogy} = teaching children; \textbf{Andragogy} = teaching adults (Knowles); \textbf{Heutagogy} = self-determined learning.
\item
  \textbf{Piaget stages}: Sensorimotor (0--2) → Pre-operational (2--7) → Concrete operational (7--11) → Formal operational (11+).
\item
  \textbf{Vygotsky}: Zone of Proximal Development (ZPD), \textbf{scaffolding}, social constructivism.
\item
  \textbf{Gardner} multiple intelligences: linguistic, logical-mathematical, spatial, musical, bodily-kinesthetic, interpersonal, intrapersonal, naturalistic (8).
\item
  \textbf{VAK} learning styles: Visual, Auditory, Kinesthetic.
\end{itemize}

\needspace{135pt}

\hypertarget{p1-01--3-factors-affecting-teaching}{%
\section{3. Factors Affecting Teaching}\label{p1-01--3-factors-affecting-teaching}}

\begingroup\small

\begin{longtable}[]{@{}
  >{\raggedright\arraybackslash}p{(\columnwidth - 2\tabcolsep) * \real{0.1825}}
  >{\raggedright\arraybackslash}p{(\columnwidth - 2\tabcolsep) * \real{0.4429}}@{}}
\toprule\noalign{}
\begin{minipage}[b]{\linewidth}\raggedright
\textbf{Factor}
\end{minipage} & \begin{minipage}[b]{\linewidth}\raggedright
\textbf{Examples}
\end{minipage} \\
\midrule\noalign{}
\endhead
\bottomrule\noalign{}
\endlastfoot
Teacher & Subject knowledge, communication, personality, motivation \\
Learner & Aptitude, interest, motivation, readiness, prior knowledge \\
Support material & Textbooks, A/V aids, e-content \\
Instructional facilities & Labs, library, smart classroom \\
Learning environment & Class size, seating, light/\allowbreak{}ventilation, classroom climate \\
Institution & Policies, administration, culture \\
\end{longtable}

\endgroup

\needdiagram{figures/p1-01-03}{55pt}

\hypertarget{p1-01--4-methods-of-teaching}{%
\section{4. Methods of Teaching}\label{p1-01--4-methods-of-teaching}}

\diagramhere{figures/p1-01-03}

\begingroup\small

\begin{longtable}[]{@{}
  >{\raggedright\arraybackslash}p{(\columnwidth - 4\tabcolsep) * \real{0.2301}}
  >{\raggedright\arraybackslash}p{(\columnwidth - 4\tabcolsep) * \real{0.3160}}
  >{\raggedright\arraybackslash}p{(\columnwidth - 4\tabcolsep) * \real{0.2733}}@{}}
\toprule\noalign{}
\begin{minipage}[b]{\linewidth}\raggedright
\textbf{Method}
\end{minipage} & \begin{minipage}[b]{\linewidth}\raggedright
\textbf{Key idea}
\end{minipage} & \begin{minipage}[b]{\linewidth}\raggedright
\textbf{Proponent}
\end{minipage} \\
\midrule\noalign{}
\endhead
\bottomrule\noalign{}
\endlastfoot
Project method & Learning by doing, purposeful activity & W.H. Kilpatrick (based on Dewey) \\
Heuristic & Student as discoverer & H.E. Armstrong \\
Programmed instruction & Small steps, immediate feedback, linear & B.F. Skinner \\
Branching programming & Branches based on response & Norman Crowder \\
Dalton plan & Contract-based individual work & Helen Parkhurst \\
Montessori & Child-centred, didactic apparatus & Maria Montessori \\
Kindergarten & "Garden of children" & Froebel \\
Basic education (Nai Talim) & Craft-centred & Gandhi \\
Flipped classroom & Content at home, practice in class & Bergmann \& Sams \\
\end{longtable}

\endgroup

\textbf{Group discussion formats}: Seminar (paper + discussion), Symposium (several speakers on aspects of one topic),
Panel discussion (informal conversation among experts before audience), Workshop (hands-on).

\needspace{135pt}

\hypertarget{p1-01--online-teaching--government-initiatives-asked-every-cycle}{%
\subsection{Online Teaching --- Government Initiatives (asked every cycle)}\label{p1-01--online-teaching--government-initiatives-asked-every-cycle}}

\begingroup\small

\begin{longtable}[]{@{}
  >{\raggedright\arraybackslash}p{(\columnwidth - 2\tabcolsep) * \real{0.2660}}
  >{\raggedright\arraybackslash}p{(\columnwidth - 2\tabcolsep) * \real{0.7142}}@{}}
\toprule\noalign{}
\begin{minipage}[b]{\linewidth}\raggedright
\textbf{Initiative}
\end{minipage} & \begin{minipage}[b]{\linewidth}\raggedright
\textbf{What it is}
\end{minipage} \\
\midrule\noalign{}
\endhead
\bottomrule\noalign{}
\endlastfoot
\textbf{SWAYAM} & Study Webs of Active-Learning for Young Aspiring Minds --- MOOC platform (2017), 4 quadrants \\
SWAYAM 4 quadrants & e-Tutorial (video), e-Content (text), Self-assessment (quiz), Discussion forum \\
\textbf{SWAYAM PRABHA} & DTH educational channels (24×7), launched with 32, now 40+ channels, via GSAT-15 \\
\textbf{NPTEL} & IITs + IISc engineering courses \\
\textbf{NDL} & National Digital Library (IIT Kharagpur) \\
\textbf{e-PG Pathshala} & PG e-content by UGC/\allowbreak{}INFLIBNET \\
\textbf{e-Yantra} & Robotics (IIT Bombay) \\
\textbf{Shodhganga} & Repository of Indian theses (INFLIBNET) \\
\textbf{Shodhgangotri} & Repository of synopses \\
\textbf{DIKSHA} & School teachers\textquotesingle{} platform (One Nation One Digital Platform) \\
\textbf{Virtual Labs} & MHRD remote labs \\
\textbf{Spoken Tutorial} & IIT Bombay software training \\
\textbf{NMEICT} & National Mission on Education through ICT \\
National coordinators of SWAYAM & AICTE, NPTEL, UGC, CEC, NCERT, NIOS, IGNOU, IIMB, NITTTR \\
\end{longtable}

\endgroup

\textbf{MOOC}: Massive Open Online Course. Term coined by Dave Cormier (2008). cMOOC (connectivist) vs xMOOC (extended, structured).

\needdiagram{figures/p1-01-04}{55pt}

\hypertarget{p1-01--5-teaching-support-system}{%
\section{5. Teaching Support System}\label{p1-01--5-teaching-support-system}}

\diagramhere{figures/p1-01-04}

\textbf{Edgar Dale\textquotesingle s Cone of Experience} (most concrete at base):

\begin{asciiblock}{142pt}
\begin{Verbatim}[fontsize=\fontsize{8.8}{10.7}\selectfont]
               /\        Verbal symbols         ← most abstract
              /  \       Visual symbols
             /    \      Recordings, radio, still pictures
            /      \     Motion pictures
           /        \    Television
          /          \   Exhibits
         /            \  Field trips
        /              \ Demonstrations
       /                \Dramatised experiences
      /                  \Contrived experiences
     /____________________\ Direct purposeful experiences ← most concrete
\end{Verbatim}
\end{asciiblock}

People remember \textasciitilde10\% of what they read, \textasciitilde90\% of what they say and do.

Audio aids: radio, tape recorder, podcast. Visual: chart, map, blackboard. \textbf{Audio-visual}: TV, film, video.
Projected aids: OHP, LCD; Non-projected: blackboard, charts, models.

\needdiagram{figures/p1-01-05}{55pt}

\hypertarget{p1-01--6-evaluation-systems}{%
\section{6. Evaluation Systems}\label{p1-01--6-evaluation-systems}}

\diagramhere{figures/p1-01-05}

\begingroup\small

\begin{longtable}[]{@{}
  >{\raggedright\arraybackslash}p{(\columnwidth - 2\tabcolsep) * \real{0.1901}}
  >{\raggedright\arraybackslash}p{(\columnwidth - 2\tabcolsep) * \real{0.5033}}@{}}
\toprule\noalign{}
\begin{minipage}[b]{\linewidth}\raggedright
\textbf{Term}
\end{minipage} & \begin{minipage}[b]{\linewidth}\raggedright
\textbf{Meaning}
\end{minipage} \\
\midrule\noalign{}
\endhead
\bottomrule\noalign{}
\endlastfoot
Measurement & Quantitative description (numbers) \\
Assessment & Collecting information about learning \\
Evaluation & Measurement + value judgement (qualitative + quantitative) \\
CCE & Continuous \& Comprehensive Evaluation --- scholastic + co-scholastic \\
Reliability & Consistency of scores \\
Validity & Measures what it intends to measure (most important) \\
Objectivity & Free from examiner bias \\
Usability/\allowbreak{}Practicability & Easy to administer \\
\end{longtable}

\endgroup

\textbf{Evaluation in CBCS}: credit = 1 hr lecture/week per semester (or 2 hrs practical). Letter grades:
O (10), A+ (9), A (8), B+ (7), B (6), C (5), P (4), F (0), Ab.

\begin{itemize}
\tightlist
\item
  \textbf{SGPA} = Σ(Cᵢ × Gᵢ) / ΣCᵢ (semester) · \textbf{CGPA} = Σ(Cᵢ × Sᵢ) / ΣCᵢ (across semesters)
\item
  Percentage ≈ CGPA × 9.5 (UGC rule-of-thumb).
\item
  CBCS course types: Core, Elective (Discipline-specific, Generic), Ability Enhancement (AECC, SEC).
\end{itemize}

\textbf{Computer-based testing \& innovations}: CAT (Computer Adaptive Testing --- difficulty adapts to answers), open-book
exams, online proctoring, e-portfolios, rubrics, peer assessment, grading instead of marks, question banks.

\needdiagram{figures/p1-01-06}{95pt}

\hypertarget{p1-01--7-deeper-dive--learning-theories-micro-teaching--teaching-models}{%
\section{7. Deeper Dive --- Learning Theories, Micro-teaching \& Teaching Models}\label{p1-01--7-deeper-dive--learning-theories-micro-teaching--teaching-models}}

\needspace{100pt}

\hypertarget{p1-01--71-learning-theories-asked-in-almost-every-cycle}{%
\subsection{7.1 Learning Theories (asked in almost every cycle)}\label{p1-01--71-learning-theories-asked-in-almost-every-cycle}}

\diagramhere{figures/p1-01-06}

\begingroup\small

\begin{longtable}[]{@{}
  >{\raggedright\arraybackslash}p{(\columnwidth - 4\tabcolsep) * \real{0.0869}}
  >{\raggedright\arraybackslash}p{(\columnwidth - 4\tabcolsep) * \real{0.4704}}
  >{\raggedright\arraybackslash}p{(\columnwidth - 4\tabcolsep) * \real{0.4427}}@{}}
\toprule\noalign{}
\begin{minipage}[b]{\linewidth}\raggedright
\textbf{Theorist}
\end{minipage} & \begin{minipage}[b]{\linewidth}\raggedright
\textbf{Key idea}
\end{minipage} & \begin{minipage}[b]{\linewidth}\raggedright
\textbf{Classroom implication}
\end{minipage} \\
\midrule\noalign{}
\endhead
\bottomrule\noalign{}
\endlastfoot
\textbf{Pavlov} & Classical conditioning: neutral stimulus + unconditioned stimulus → conditioned response (dog \& bell) & Pleasant classroom associations reduce fear of a subject \\
\textbf{Thorndike} & Connectionism; laws of \textbf{readiness, exercise, effect} & Practice, rewards, readiness before teaching \\
\textbf{Skinner} & Operant conditioning; positive/\allowbreak{}negative reinforcement, punishment, shaping, schedules of reinforcement & Programmed instruction, immediate feedback \\
\textbf{Kohler} & Insight learning (chimpanzee Sultan) --- sudden "aha" grasp of relationships & Present whole problems \\
\textbf{Bruner} & Discovery learning; modes: \textbf{enactive → iconic → symbolic}; spiral curriculum & Revisit topics with increasing depth \\
\textbf{Ausubel} & Meaningful verbal learning; \textbf{advance organisers} & Start a lesson with an overview linking to prior knowledge \\
\textbf{Bandura} & Social learning / observational learning (Bobo doll); modelling; self-efficacy & Teacher as role model \\
\textbf{Gagné} & Conditions of learning; 8 types of learning; \textbf{9 events of instruction} & Structured lesson design \\
\textbf{Maslow} & Physiological → safety → love/\allowbreak{}belonging → esteem → self-actualisation & Basic needs before academic growth \\
\end{longtable}

\endgroup

\begin{itemize}
\tightlist
\item
  \textbf{Negative reinforcement} = removing an unpleasant stimulus to \emph{increase} behaviour (not punishment).
\item
  Gagné\textquotesingle s 9 events: gain attention → inform objectives → recall prior learning → present content → provide guidance → elicit performance → give feedback → assess performance → enhance retention and transfer.
\end{itemize}

\needdiagram{figures/p1-01-07}{90pt}

\hypertarget{p1-01--72-micro-teaching}{%
\subsection{7.2 Micro-teaching}\label{p1-01--72-micro-teaching}}

Developed at \textbf{Stanford University (1961--63) by Dwight Allen} and colleagues. A scaled-down teaching encounter: one skill, 5--10 students, 5--10 minutes.

\diagramhere{figures/p1-01-07}

\begin{itemize}
\tightlist
\item
  Indian model (NCERT/Passi): total cycle \textbf{36 minutes}.
\item
  Skills: introducing a lesson, questioning (fluency, probing), explaining, illustrating with examples, \textbf{reinforcement}, stimulus variation, blackboard writing, closure, using A/V aids.
\item
  \textbf{Simulated teaching} = role play of teaching in a training situation; \textbf{team teaching} = two or more teachers plan and teach together.
\end{itemize}

\needspace{135pt}

\hypertarget{p1-01--73-models-of-teaching}{%
\subsection{7.3 Models of Teaching}\label{p1-01--73-models-of-teaching}}

\begingroup\small

\begin{longtable}[]{@{}
  >{\raggedright\arraybackslash}p{(\columnwidth - 4\tabcolsep) * \real{0.1767}}
  >{\raggedright\arraybackslash}p{(\columnwidth - 4\tabcolsep) * \real{0.1777}}
  >{\raggedright\arraybackslash}p{(\columnwidth - 4\tabcolsep) * \real{0.6456}}@{}}
\toprule\noalign{}
\begin{minipage}[b]{\linewidth}\raggedright
\textbf{Model}
\end{minipage} & \begin{minipage}[b]{\linewidth}\raggedright
\textbf{Proponent}
\end{minipage} & \begin{minipage}[b]{\linewidth}\raggedright
\textbf{Key steps / idea}
\end{minipage} \\
\midrule\noalign{}
\endhead
\bottomrule\noalign{}
\endlastfoot
\textbf{Basic Teaching Model} & Robert Glaser (1962) & Instructional objectives → Entering behaviour → Instructional procedures → Performance assessment (with feedback loop) \\
Concept Attainment & Jerome Bruner & Examples \& non-examples → identify concept attributes \\
Inductive Thinking & Hilda Taba & Concept formation → interpretation of data → application \\
Advance Organiser & David Ausubel & Organiser presented before content \\
Inquiry Training & Richard Suchman & Puzzling event → students ask yes/no questions → form theories \\
Jurisprudential & Oliver \& Shaver & Debating public issues \\
Synectics & William Gordon & Creativity through metaphors \& analogies \\
\end{longtable}

\endgroup

Families of teaching models (\textbf{Joyce \& Weil}): Information processing, Social interaction, Personal, Behavioural systems.

\begin{asciiblock}{89pt}
\begin{Verbatim}[fontsize=\fontsize{8.8}{10.7}\selectfont]
Glaser's Basic Teaching Model
 ┌──────────────┐   ┌──────────────┐   ┌──────────────┐   ┌──────────────┐
 │ Instructional├──►│  Entering    ├──►│ Instructional├──►│ Performance  │
 │ objectives   │   │  behaviour   │   │ procedures   │   │ assessment   │
 └──────▲───────┘   └──────▲───────┘   └──────▲───────┘   └──────┬───────┘
        └──────────────────┴──────────────────┴── feedback ───────┘
\end{Verbatim}
\end{asciiblock}

\needspace{100pt}

\hypertarget{p1-01--74-questioning--classroom-management}{%
\subsection{7.4 Questioning \& Classroom Management}\label{p1-01--74-questioning--classroom-management}}

\begin{itemize}
\tightlist
\item
  \textbf{Convergent} questions (one correct answer, lower order) vs \textbf{divergent} (many answers, creativity).
\item
  \textbf{Wait time} (Mary Budd Rowe): waiting 3--5 seconds after a question improves answer quality.
\item
  \textbf{Kounin}: "withitness" (teacher aware of everything), overlapping, momentum, smoothness.
\item
  Teacher leadership styles (Lewin): \textbf{autocratic, democratic, laissez-faire} --- democratic is best for learning climate.
\item
  \textbf{Hidden curriculum}: values and norms learnt implicitly from school culture.
\end{itemize}

\needspace{135pt}

\hypertarget{p1-01--previous-year-questions-pyq-pattern}{%
\section{Previous Year Questions (PYQ pattern)}\label{p1-01--previous-year-questions-pyq-pattern}}

\textbf{Levels of teaching}

\begin{enumerate}
\def\labelenumi{\arabic{enumi}.}
\tightlist
\item
  The "reflective level" of teaching is associated with: (a) Herbart (b) Morrison (c) Hunt (d) Skinner --- \textbf{Ans: (c)}
\item
  Which level of teaching is the most "thoughtless"? (a) Memory (b) Understanding (c) Reflective (d) Autonomous --- \textbf{Ans: (a)}
\item
  Arrange the levels from lowest to highest: \textbf{Memory → Understanding → Reflective}.
\end{enumerate}

\textbf{Bloom\textquotesingle s taxonomy}

\begin{enumerate}
\def\labelenumi{\arabic{enumi}.}
\setcounter{enumi}{3}
\tightlist
\item
  In revised Bloom\textquotesingle s taxonomy the highest cognitive level is: (a) Evaluating (b) Creating (c) Analysing (d) Synthesis --- \textbf{Ans: (b)}
\item
  "Characterisation by a value" belongs to which domain? --- \textbf{Ans: Affective (Krathwohl)}
\item
  Match: Remembering--recall; Applying--use; Analysing--differentiate; Evaluating--judge. (Standard match question.)
\end{enumerate}

\textbf{Learner characteristics}

\begin{enumerate}
\def\labelenumi{\arabic{enumi}.}
\setcounter{enumi}{6}
\tightlist
\item
  Andragogy refers to: (a) teaching children (b) teaching adults (c) self-learning (d) group learning --- \textbf{Ans: (b)}
\item
  ZPD is a concept given by: \textbf{Vygotsky}
\item
  Which is NOT a characteristic of adult learners? (a) self-directed (b) experience-based (c) externally driven (d) problem-centred --- \textbf{Ans: (c)}
\end{enumerate}

\textbf{Methods}

\begin{enumerate}
\def\labelenumi{\arabic{enumi}.}
\setcounter{enumi}{9}
\tightlist
\item
  Programmed instruction is based on the theory of: (a) Classical conditioning (b) Operant conditioning (c) Insight (d) Constructivism --- \textbf{Ans: (b) Skinner}
\item
  "Learning by doing" is the core of: \textbf{Project method (Kilpatrick)}
\item
  Which is a learner-centred method? (a) Lecture (b) Demonstration (c) Problem solving (d) Team teaching --- \textbf{Ans: (c)}
\item
  In a flipped classroom, students first encounter new content: \textbf{at home (videos/readings)}
\end{enumerate}

\textbf{Online initiatives}

\begin{enumerate}
\def\labelenumi{\arabic{enumi}.}
\setcounter{enumi}{13}
\tightlist
\item
  SWAYAM PRABHA provides: (a) MOOCs (b) 24×7 DTH educational channels (c) Digital library (d) Thesis repository --- \textbf{Ans: (b)}
\item
  How many quadrants does a SWAYAM course have? --- \textbf{Ans: 4}
\item
  Shodhganga is maintained by: \textbf{INFLIBNET (Gandhinagar)}
\item
  The national coordinator for engineering courses on SWAYAM: \textbf{NPTEL}
\end{enumerate}

\textbf{Evaluation}

\begin{enumerate}
\def\labelenumi{\arabic{enumi}.}
\setcounter{enumi}{17}
\tightlist
\item
  Evaluation that takes place during the teaching-learning process is: \textbf{Formative}
\item
  A test that measures what it intends to measure is: \textbf{Valid}
\item
  Consistency of test scores refers to: \textbf{Reliability}
\item
  Norm-referenced test compares a student with: \textbf{performance of peer group}
\item
  Diagnostic evaluation is meant to: \textbf{identify learning difficulties}
\item
  In CBCS, the grade point for "O" (Outstanding) is: \textbf{10}
\end{enumerate}

\textbf{Teaching aids}

\begin{enumerate}
\def\labelenumi{\arabic{enumi}.}
\setcounter{enumi}{23}
\tightlist
\item
  In Edgar Dale\textquotesingle s Cone, the most concrete experience is: \textbf{Direct purposeful experience}
\item
  Which is a non-projected aid? (a) LCD (b) OHP (c) Blackboard (d) Slide --- \textbf{Ans: (c)}
\end{enumerate}

\textbf{Statement-based (common format)}

\begin{enumerate}
\def\labelenumi{\arabic{enumi}.}
\setcounter{enumi}{25}
\tightlist
\item
  Statement I: Effective teaching is learner-centred. Statement II: The teacher\textquotesingle s role is only to transmit content.
  --- \textbf{I true, II false}
\end{enumerate}

\textbf{More practice questions}

\begin{enumerate}
\def\labelenumi{\arabic{enumi}.}
\setcounter{enumi}{26}
\tightlist
\item
  Learning by insight was demonstrated by: \textbf{Wolfgang Kohler}
\item
  "Law of effect" is associated with: \textbf{E.L. Thorndike}
\item
  Advance organisers were proposed by: \textbf{David Ausubel}
\item
  Bruner\textquotesingle s three modes of representation in order: \textbf{Enactive, Iconic, Symbolic}
\item
  Removing an unpleasant stimulus to increase a desired behaviour is: \textbf{negative reinforcement}
\item
  Micro-teaching was developed at: \textbf{Stanford University}
\item
  Duration of one micro-teaching cycle in the Indian model: \textbf{36 minutes}
\item
  The basic teaching model was given by: \textbf{Robert Glaser}
\item
  The first component of Glaser\textquotesingle s model is: \textbf{instructional objectives}
\item
  Concept attainment model is associated with: \textbf{Jerome Bruner}
\item
  A question with several acceptable answers is: \textbf{divergent}
\item
  Observational (social) learning theory was given by: \textbf{Albert Bandura}
\item
  In Maslow\textquotesingle s hierarchy, the highest need is: \textbf{self-actualisation}
\item
  Which leadership style creates the most positive classroom climate? \textbf{Democratic}
\item
  Statement I: Formative evaluation is diagnostic in nature. Statement II: Summative evaluation is used for certification. --- \textbf{Both true}
\item
  Match: Kilpatrick--Project method; Skinner--Programmed learning; Armstrong--Heuristic; Parkhurst--Dalton plan.
\item
  Which is NOT a characteristic of effective teaching? (a) learner involvement (b) feedback (c) rote memorisation (d) use of examples --- \textbf{Ans: (c)}
\end{enumerate}

\begin{revbox}

\begin{itemize}
\tightlist
\item
  Memory--Herbart · Understanding--Morrison · Reflective--Hunt
\item
  Formative = during; Summative = end; Diagnostic = difficulties; Placement = before
\item
  Validity \textgreater{} Reliability (valid test is always reliable, not vice versa)
\item
  SWAYAM (2017, MOOCs, 4 quadrants) · SWAYAM PRABHA (DTH channels, 32 at launch → 40+) · Shodhganga (theses, INFLIBNET)
\end{itemize}

\end{revbox}
